\documentclass[10pt,a4paper]{amsart}
\usepackage{amsmath,amssymb,amsthm,graphicx}
% Retained minimal preamble for the Stacks Project chapter sites-cohomology.tex (Cohomology on Sites).
% Verbatim excerpt of the chapter preamble: Xy-pic setup (native preamble lines 14-18) and the native
% theorem declarations with the mathematical macros (native preamble lines 154-217). Copied mechanically;
% nothing added, removed or redefined. The wrapper class and the amsmath/amssymb/amsthm/graphicx base
% packages are supplied by the assembler, so the class/externaldocument part of the native preamble is omitted.
% For commutative diagrams we use Xy-pic
\usepackage[all]{xy}

% We use 2cell for 2-commutative diagrams.
\xyoption{2cell}
% Theorem environments.
%
\theoremstyle{plain}
\newtheorem{theorem}[subsection]{Theorem}
\newtheorem{proposition}[subsection]{Proposition}
\newtheorem{lemma}[subsection]{Lemma}

\theoremstyle{definition}
\newtheorem{definition}[subsection]{Definition}
\newtheorem{example}[subsection]{Example}
\newtheorem{exercise}[subsection]{Exercise}
\newtheorem{situation}[subsection]{Situation}

\theoremstyle{remark}
\newtheorem{remark}[subsection]{Remark}
\newtheorem{remarks}[subsection]{Remarks}

\numberwithin{equation}{subsection}

% Macros
%
\def\lim{\mathop{\mathrm{lim}}\nolimits}
\def\colim{\mathop{\mathrm{colim}}\nolimits}
\def\Spec{\mathop{\mathrm{Spec}}}
\def\Hom{\mathop{\mathrm{Hom}}\nolimits}
\def\Ext{\mathop{\mathrm{Ext}}\nolimits}
\def\SheafHom{\mathop{\mathcal{H}\!\mathit{om}}\nolimits}
\def\SheafExt{\mathop{\mathcal{E}\!\mathit{xt}}\nolimits}
\def\Sch{\mathit{Sch}}
\def\Mor{\mathop{\mathrm{Mor}}\nolimits}
\def\Ob{\mathop{\mathrm{Ob}}\nolimits}
\def\Sh{\mathop{\mathit{Sh}}\nolimits}
\def\NL{\mathop{N\!L}\nolimits}
\def\CH{\mathop{\mathrm{CH}}\nolimits}
\def\proetale{{pro\text{-}\acute{e}tale}}
\def\etale{{\acute{e}tale}}
\def\QCoh{\mathit{QCoh}}
\def\Ker{\mathop{\mathrm{Ker}}}
\def\Im{\mathop{\mathrm{Im}}}
\def\Coker{\mathop{\mathrm{Coker}}}
\def\Coim{\mathop{\mathrm{Coim}}}

% Boxtimes
%
\DeclareMathSymbol{\boxtimes}{\mathbin}{AMSa}{"02}

%
% Macros for moduli stacks/spaces
%
\def\QCohstack{\mathcal{QC}\!\mathit{oh}}
\def\Cohstack{\mathcal{C}\!\mathit{oh}}
\def\Spacesstack{\mathcal{S}\!\mathit{paces}}
\def\Quotfunctor{\mathrm{Quot}}
\def\Hilbfunctor{\mathrm{Hilb}}
\def\Curvesstack{\mathcal{C}\!\mathit{urves}}
\def\Polarizedstack{\mathcal{P}\!\mathit{olarized}}
\def\Complexesstack{\mathcal{C}\!\mathit{omplexes}}
% \Pic is the operator that assigns to X its picard group, usage \Pic(X)
% \Picardstack_{X/B} denotes the Picard stack of X over B
% \Picardfunctor_{X/B} denotes the Picard functor of X over B
\def\Pic{\mathop{\mathrm{Pic}}\nolimits}
\def\Picardstack{\mathcal{P}\!\mathit{ic}}
\def\Picardfunctor{\mathrm{Pic}}
\def\Deformationcategory{\mathcal{D}\!\mathit{ef}}

% --- appended 2026-10-12 (Chief): the project own definitions of the reference/slogan/history
% display environments, copied mechanically from the shared Stacks preamble lines 8-12
% (\usepackage{verbatim} plus the three \newenvironment lines, which the chapter preamble keeps
% verbatim from the project preamble). Nothing added, removed or redefined: these environments
% carry no mathematical content and the project itself typesets them as comments. Without these
% three lines a counted statement whose own text uses them fails with Environment slogan undefined.
\usepackage{verbatim}
\newenvironment{reference}{\comment}{\endcomment}
\newenvironment{slogan}{\comment}{\endcomment}
\newenvironment{history}{\comment}{\endcomment}
\usepackage{cleveref}
\crefname{originalreference}{original reference}{original references}
\Crefname{originalreference}{Original reference}{Original references}
\makeatletter
\newlabel{item-A-and-P}{{0EZ7}{source}{}{sourcepointer-1}{}}
\newlabel{item-A-and-P@cref}{{[originalreference][0][]0EZ7}{[1][0][]source}{}{}{}}
\makeatother
\crefname{originalreference}{original reference}{original references}
\Crefname{originalreference}{Original reference}{Original references}
\makeatletter
\newlabel{item-A-sheaf}{{0EZ6}{source}{}{sourcepointer-2}{}}
\newlabel{item-A-sheaf@cref}{{[originalreference][0][]0EZ6}{[1][0][]source}{}{}{}}
\makeatother
\crefname{originalreference}{original reference}{original references}
\Crefname{originalreference}{Original reference}{Original references}
\makeatletter
\newlabel{item-base-change-P}{{0EZ4}{source}{}{sourcepointer-3}{}}
\newlabel{item-base-change-P@cref}{{[originalreference][0][]0EZ4}{[1][0][]source}{}{}{}}
\makeatother
\crefname{originalreference}{original reference}{original references}
\Crefname{originalreference}{Original reference}{Original references}
\makeatletter
\newlabel{item-restriction-A}{{0EZ5}{source}{}{sourcepointer-4}{}}
\newlabel{item-restriction-A@cref}{{[originalreference][0][]0EZ5}{[1][0][]source}{}{}{}}
\makeatother
\crefname{originalreference}{original reference}{original references}
\Crefname{originalreference}{Original reference}{Original references}
\makeatletter
\newlabel{lemma-cohomology-of-open}{{01E1}{source}{}{sourcepointer-5}{}}
\newlabel{lemma-cohomology-of-open@cref}{{[originalreference][0][]01E1}{[1][0][]source}{}{}{}}
\makeatother
\crefname{originalreference}{original reference}{original references}
\Crefname{originalreference}{Original reference}{Original references}
\makeatletter
\newlabel{lemma-localize-cartesian-square}{{0EYV}{source}{}{sourcepointer-6}{}}
\newlabel{lemma-localize-cartesian-square@cref}{{[originalreference][0][]0EYV}{[1][0][]source}{}{}{}}
\makeatother
\crefname{originalreference}{original reference}{original references}
\Crefname{originalreference}{Original reference}{Original references}
\makeatletter
\newlabel{situation-compare}{{0EZ3}{source}{}{sourcepointer-7}{}}
\newlabel{situation-compare@cref}{{[originalreference][0][]0EZ3}{[1][0][]source}{}{}{}}
\makeatother
\title{A class in the equalizer on the double cover becomes zero after restricting to a covering}
\author{The Stacks Project Authors}
\date{stacks-project:sites-cohomology.tex:local-snapshot-20261011; GNU FDL 1.2 or later}
\begin{document}
\maketitle
\noindent\textit{Editorial note.} All delivered mathematical text is the author's own, copied from the pinned Stacks Project chapter file sites-cohomology.tex; no statement or proof step was rewritten, reordered or removed.\par\medskip
\noindent\textit{Editorial note.} The counted proof is the author's complete argument: the proof is written out in full: the reformulation through the cartesian square, the trivial base change identification, the equalizer argument, and the closing sheaf-property computation using the four standing properties.\par\medskip
\noindent\textit{Editorial note.} The label lemma-make-class-zero-general names the counted statement, and the proof environment immediately following it in the source is delivered verbatim as its proof.\par\medskip
\noindent\textit{Editorial note.} The counted statement is carried by the author's own lemma environment in the chapter source.\par\medskip
\noindent\textit{Editorial note.} Pre-open gate: python operations/agent-authoring/tools/precard\_gate.py --source sites-cohomology.tex --label lemma-make-class-zero-general --start 7089 --end 7145 --local-id 3551 returned RESULT=OK (L3551 free, occupancy FREE). Candidate taken from lane-s4/halfB-priority.json.\par\medskip
\noindent\textit{Editorial note.} The statement and proof are copied from the pinned chapter file; the cited inputs are kept as external pointers with their canonical Stacks tags.\par\medskip
\noindent\textit{Allowed standard prerequisites (editorial).}\begin{itemize}
\item The standing comparison situation and its five properties. Cohomology on Sites, the situation comparing two topologies, together with the properties (item-base-change-P), (item-restriction-A), (item-A-and-P) and (item-A-sheaf). Supplies the existence of fibre products, the membership of the restrictions in the coefficient categories, and the sheaf property used at the end.
\item Trivial base change for the cartesian square of f. Same chapter, the lemma computing the pullback of the higher pushforward along the cartesian square of a morphism. Supplies the identification of f\textasciicircum{}\{-1\}R\textasciicircum{}\{n+1\}f\_* of the restriction with the higher pushforward of the restriction to the double cover.
\item Cohomology over an object as cohomology of the localized site. Same chapter, the lemma identifying the cohomology over an object with cohomology of the localized site. Supplies the rewriting of the cohomology groups in the statement.
\end{itemize}\medskip
\ifdefined\hypertarget\hypertarget{sourcepointer-1}{}\fi
\noindent\textit{Source reference.} Original reference 0EZ7: sources/stacks-sites-cohomology/pinned-original-source.tex, source lines 6862--6862. Prerequisite pointer invoked at proof line(s) 7139; it names a standard fact of the same chapter that the counted argument uses, and every step of the counted proof built on it is retained verbatim in the delivered body.\par\medskip
\ifdefined\hypertarget\hypertarget{sourcepointer-2}{}\fi
\noindent\textit{Source reference.} Original reference 0EZ6: sources/stacks-sites-cohomology/pinned-original-source.tex, source lines 6857--6857. Prerequisite pointer invoked at proof line(s) 7141; it names a standard fact of the same chapter that the counted argument uses, and every step of the counted proof built on it is retained verbatim in the delivered body.\par\medskip
\ifdefined\hypertarget\hypertarget{sourcepointer-3}{}\fi
\noindent\textit{Source reference.} Original reference 0EZ4: sources/stacks-sites-cohomology/pinned-original-source.tex, source lines 6849--6849. Prerequisite pointer invoked at proof line(s) 7110; it names a standard fact of the same chapter that the counted argument uses, and every step of the counted proof built on it is retained verbatim in the delivered body.\par\medskip
\ifdefined\hypertarget\hypertarget{sourcepointer-4}{}\fi
\noindent\textit{Source reference.} Original reference 0EZ5: sources/stacks-sites-cohomology/pinned-original-source.tex, source lines 6853--6853. Prerequisite pointer invoked at proof line(s) 7137; it names a standard fact of the same chapter that the counted argument uses, and every step of the counted proof built on it is retained verbatim in the delivered body.\par\medskip
\ifdefined\hypertarget\hypertarget{sourcepointer-5}{}\fi
\noindent\textit{Source reference.} Original reference 01E1: sources/stacks-sites-cohomology/pinned-original-source.tex, source lines 608--608. Prerequisite pointer invoked at proof line(s) 7115; it names a standard fact of the same chapter that the counted argument uses, and every step of the counted proof built on it is retained verbatim in the delivered body.\par\medskip
\ifdefined\hypertarget\hypertarget{sourcepointer-6}{}\fi
\noindent\textit{Source reference.} Original reference 0EYV: sources/stacks-sites-cohomology/pinned-original-source.tex, source lines 4652--4652. Prerequisite pointer invoked at proof line(s) 7126; it names a standard fact of the same chapter that the counted argument uses, and every step of the counted proof built on it is retained verbatim in the delivered body.\par\medskip
\ifdefined\hypertarget\hypertarget{sourcepointer-7}{}\fi
\noindent\textit{Source reference.} Original reference 0EZ3: sources/stacks-sites-cohomology/pinned-original-source.tex, source lines 6841--6841. Prerequisite pointer invoked at proof line(s) 7091; it names a standard fact of the same chapter that the counted argument uses, and every step of the counted proof built on it is retained verbatim in the delivered body.\par\medskip
\begin{lemma}
\label{lemma-make-class-zero-general}
In Situation \ref{situation-compare} let $f : X \to Y$
be in $\mathcal{P}$ such that $\{X \to Y\}$ is a $\tau$-covering.
Let $\mathcal{F}'$ be in $\mathcal{A}'_Y$. If $n \geq 0$ and
$$
\theta \in
\text{Equalizer}\left(
\xymatrix{
H^{n + 1}_{\tau'}(X, \mathcal{F}')
\ar@<1ex>[r] \ar@<-1ex>[r] &
H^{n + 1}_{\tau'}(X \times_Y X, \mathcal{F}')
}
\right)
$$
then there exists a $\tau'$-covering $\{Y_i \to Y\}$
such that $\theta$ restricts to zero in
$H^{n + 1}_{\tau'}(Y_i \times_Y X, \mathcal{F}')$.
\end{lemma}

\begin{proof}
Observe that $X \times_Y X$ exists by (\ref{item-base-change-P}).
For $Z$ in $\mathcal{C}/Y$ denote $\mathcal{F}'|_Z$
the restriction of $\mathcal{F}'$ to $\mathcal{C}_{\tau'}/Z$.
Recall that $H^{n + 1}_{\tau'}(X, \mathcal{F}') =
H^{n + 1}(\mathcal{C}_{\tau'}/X, \mathcal{F}'|_X)$, see
Lemma \ref{lemma-cohomology-of-open}.
The lemma asserts that the
image $\overline{\theta} \in H^0(Y, R^{n + 1}f_{\tau', *}\mathcal{F}'|_X)$
of $\theta$ is zero. Consider the cartesian diagram
$$
\xymatrix{
X \times_Y X \ar[d]_{\text{pr}_1} \ar[r]_{\text{pr}_2} &
X \ar[d]^f \\
X \ar[r]^f & Y
}
$$
By trivial base change (Lemma \ref{lemma-localize-cartesian-square})
we have
$$
f_{\tau'}^{-1}R^{n + 1}f_{\tau', *}(\mathcal{F}'|_X) =
R^{n + 1}\text{pr}_{1, \tau', *}(\mathcal{F}'|_{X \times_Y X})
$$
If $\text{pr}_1^{-1}\theta = \text{pr}_2^{-1}\theta$,
then the section $f_{\tau'}^{-1}\overline{\theta}$ of
$f_{\tau'}^{-1}R^{n + 1}f_{\tau', *}(\mathcal{F}'|_X)$ is zero,
because it is clear that $\text{pr}_1^{-1}\theta$ maps to the zero element in
$H^0(X, R^{n + 1}\text{pr}_{1, \tau', *}(\mathcal{F}'|_{X \times_Y X}))$.
By (\ref{item-restriction-A}) we have $\mathcal{F}'|_X$ in $\mathcal{A}'_X$.
Thus $\mathcal{G}' = R^{n + 1}f_{\tau', *}(\mathcal{F}'|_X)$
is an object of $\mathcal{A}'_Y$ by (\ref{item-A-and-P}).
Thus $\mathcal{G}'$ satisfies the sheaf property for
$\tau$-coverings by (\ref{item-A-sheaf}).
Since $\{X \to Y\}$ is a $\tau$-covering we conclude
that restriction $\mathcal{G}'(Y) \to \mathcal{G}'(X)$ is injective.
It follows that $\overline{\theta}$ is zero.
\end{proof}

\end{document}
